\begin{pastexam}[label={px:nchu-111-is-07}]{中興資工所乙組 111 學年度 第 7 題}
（問答）

(7 pts) Draw the contents of the two open addressing hash tables in the boxes below. The size of the hash
table is 9. The hash function used is H(k) = k mod 9. What values will be in the hash table after the following
sequence of insertions? Draw the values in the boxes below, and show your work for partial credit.
\begin{center}
18, 16, 10, 7, 26
\end{center}
\begin{center}
\begin{tabular}{c@{\hspace{3cm}}c}
(a). Linear Probing & (b). Quadratic Probing\\[2pt]
\begin{tabular}{c|p{1.2cm}|}
\cline{2-2}
0 & \\ \cline{2-2}
1 & \\ \cline{2-2}
2 & \\ \cline{2-2}
3 & \\ \cline{2-2}
4 & \\ \cline{2-2}
5 & \\ \cline{2-2}
6 & \\ \cline{2-2}
7 & \\ \cline{2-2}
8 & \\ \cline{2-2}
\end{tabular}
&
\begin{tabular}{c|p{1.2cm}|}
\cline{2-2}
0 & \\ \cline{2-2}
1 & \\ \cline{2-2}
2 & \\ \cline{2-2}
3 & \\ \cline{2-2}
4 & \\ \cline{2-2}
5 & \\ \cline{2-2}
6 & \\ \cline{2-2}
7 & \\ \cline{2-2}
8 & \\ \cline{2-2}
\end{tabular}
\end{tabular}
\end{center}
\end{pastexam}

\begin{pastsoln}
題目沒寫探測公式，照資結書第~12.3~節：線性探測試 $H(k) + i$，二次探測試 $H(k) + i^2$，
$i = 0, 1, 2, \ldots$，都取 $\bmod 9$。

起始位置：$18 \bmod 9 = 0$，$16 \bmod 9 = 7$，$10 \bmod 9 = 1$，$7 \bmod 9 = 7$，$26 \bmod 9 = 8$。
前三個鍵都直接放進去，衝突從 $7$ 開始。

\textbf{(a) 線性探測}。
\begin{itemize}
\item $18 \to 0$，$16 \to 7$，$10 \to 1$。
\item $7$：第 $7$ 格有 $16$，試第 $8$ 格，空的，放入。
\item $26$：第 $8$ 格有 $7$，往後試第 $0$ 格（有 $18$）、第 $1$ 格（有 $10$）、第 $2$ 格，空的，放入。
\end{itemize}

\textbf{(b) 二次探測}。
\begin{itemize}
\item $18 \to 0$，$16 \to 7$，$10 \to 1$。
\item $7$：$7 + 0 = 7$ 有 $16$；$7 + 1 = 8$，空的，放入。
\item $26$：$8 + 0 = 8$ 有 $7$；$8 + 1 = 9 \bmod 9 = 0$ 有 $18$；$8 + 4 = 12 \bmod 9 = 3$，空的，放入。
\end{itemize}

結果：
\begin{center}
\begin{tabular}{c@{\hspace{2.5cm}}c}
(a) 線性探測 & (b) 二次探測\\[2pt]
\begin{tabular}{c|c|}
\cline{2-2}
0 & 18 \\ \cline{2-2}
1 & 10 \\ \cline{2-2}
2 & 26 \\ \cline{2-2}
3 & \\ \cline{2-2}
4 & \\ \cline{2-2}
5 & \\ \cline{2-2}
6 & \\ \cline{2-2}
7 & 16 \\ \cline{2-2}
8 & 7 \\ \cline{2-2}
\end{tabular}
&
\begin{tabular}{c|c|}
\cline{2-2}
0 & 18 \\ \cline{2-2}
1 & 10 \\ \cline{2-2}
2 & \\ \cline{2-2}
3 & 26 \\ \cline{2-2}
4 & \\ \cline{2-2}
5 & \\ \cline{2-2}
6 & \\ \cline{2-2}
7 & 16 \\ \cline{2-2}
8 & 7 \\ \cline{2-2}
\end{tabular}
\end{tabular}
\end{center}
線性探測時 $26$ 要經過 $0, 1$ 兩個被佔的格子，這就是主叢集（primary clustering）：$0$--$1$ 和 $7$--$8$
兩段連在一起（表是環狀的），落在其中任何一格的鍵都要一路走到尾端。 二次探測跳得比較開，$26$ 只試了
三格。
\end{pastsoln}
